\section{How to Start and Win the Game}\label{sec:15.0}

\ruleslabel{General Rule}
\begin{generalrule}
Before the start of play, both Players must deploy a number of units on the map
according to the instructions of Cases 15.1 and 15.2. In addition, the Axis
Player must calculate the number of Anti-Guerrilla Operations that he may
perform during the course of the game (see 15.3).
\end{generalrule}

\ruleslabel{Cases}

\subsection{Axis Set-Up}\label{sec:15.1}

Units to be deployed are indicated by nationality, designation, and strength
(in parentheses). Axis units may be deployed in any Objective Display of the
indicated Zone.

\begin{description}[style=nextline,leftmargin=1em,labelindent=0pt,font=\bfseries]
  \item[A. German]
  \begin{enumerate}
    \item Serbia: 704(12), 714(12), 717(12).
    \item Bosnia: 718(12).
  \end{enumerate}

  \item[B. Italian]
  \begin{enumerate}
    \item Slovenia: Emilia(6), Marche(6), Alpi Gr(6), CD Alpi(6).
    \item Istria: Murge(6), Messina(6).
    \item Croatia: Bergamo(6), Sassari(6).
    \item Dalmatia: Savoia(6), Zara(6).
    \item Montenegro: Re(6), Lombard(6), Macerata(6), Isonzo(6), Ferrara(6),
      Taurinense(6), Perugia(6), Venezia(6), Puglia(6).
    \item Albania: Firenze(6), Arezzo(6), Parma(6).
  \end{enumerate}

  \item[C. Serbian]
  \begin{enumerate}
    \item Serbia: 1(1), 2(1), 3(1), 4(1), 5(1).
  \end{enumerate}

  \item[D. Bulgarian]
  \begin{enumerate}
    \item Macedonia: 14(12), 22(12), 24(12).
  \end{enumerate}
\end{description}

\subsection{Yugoslav Set-Up}\label{sec:15.2}

The Yugoslav Player deploys 5 Chetnik groups in the Yugoslav/Chetnik Hide-away
circle of the Serbia Zone Display and 2 Chetnik groups in the
Yugoslav/Chetnik Hide-away circle of the Montenegro Zone Display. These units
are under the control of the Yugoslav Player at the start of the game.

\subsection{Axis Anti-Guerrilla Limitations}\label{sec:15.3}

\case{15.31} Before starting play, the Axis Player should place all ten Soviet
units available in the counter mix in a cup. Each Soviet unit is back-printed
with a number of Anti-Guerrilla Operations (``AGO''). The Axis Player blindly
chooses one Soviet unit. The number on its reverse is the number of
Anti-Guerrilla Operations the Axis Player may perform during the game.

\case{15.32} The chosen Soviet counter is kept hidden in front of the Axis
Player between Game-Turns 1 and 14 (inclusive). Each time an Anti-Guerrilla
Operation is performed, the Axis Player records it on a separate sheet of
paper. At the beginning of Game-Turn 15, the counter is revealed to the
Yugoslav Player, who verifies that the number of Anti-Guerrilla Operations
performed does not exceed the counter's number. During this Game-Turn, all
Soviet units become available as reinforcements.

\subsection{Special Rules (Game-Turns One and Two Only)}\label{sec:15.4}

On Game-Turns 1 and 2, all of both Players' units must remain in the Zone in
which they start the game (or enter the game as reinforcements).

\subsection{How to Win the Game}\label{sec:15.5}

The game ends at the completion of the Victory Point Stage of Game-Turn 17 (the
remaining Phases of this Game-Turn are deleted). At this time, the final
Yugoslav Victory Point total is computed and a winner is determined according
to the following schedule:

\begin{center}
\begin{tabularx}{\linewidth}{@{}Xr@{}}
  \toprule
  Result & Yugoslav Victory Points \\
  \midrule
  Yugoslav Decisive Victory & Greater than 800 \\
  Yugoslav Substantive Victory & 701--800 \\
  Yugoslav Marginal Victory & 601--700 \\
  Axis Marginal Victory & 501--600 \\
  Axis Substantive Victory & 401--500 \\
  Axis Decisive Victory & Less than 401 \\
  \bottomrule
\end{tabularx}
\end{center}
